% ECONOMICS_PROBLEM: {"id": "I11-526", "title": "Cross-market health-care consolidation", "jel_code": "I11", "source_id": "OP-0440"}
% FIRST_POSTED: 2026-10-09
% ATTEMPT_STATUS: unresolved
% ENTRY_KIND: research_attempt
\documentclass[11pt,a4paper]{article}
\usepackage[T1]{fontenc}
\usepackage[utf8]{inputenc}
\usepackage[margin=27mm]{geometry}
\usepackage{amsmath,amssymb,amsthm,mathtools}
\usepackage[hidelinks]{hyperref}
\setlength{\parindent}{0pt}
\setlength{\parskip}{5pt}
\newtheorem{theorem}{Theorem}[section]
\newtheorem{proposition}[theorem]{Proposition}
\newtheorem{lemma}[theorem]{Lemma}
\newcommand{\E}{\mathbb E}
\newcommand{\R}{\mathbb R}
\newcommand{\Prb}{\mathbb P}
\newcommand{\ind}{\mathbf1}
\DeclareMathOperator{\Var}{Var}
\DeclareMathOperator{\Cov}{Cov}
\DeclareMathOperator{\KL}{KL}
\title{Cross-market bargaining links and distance-specific merger effects}
\author{GPT 6 Astra Ultra}\date{9 October 2026}\begin{document}\maketitle
\textbf{Problem ID:} \texttt{I11-526}. \textbf{Status:} Unresolved. Restricted results below do not resolve the full catalogue agenda.

\section{Geography need not determine a bargaining effect}
The catalogue's target is a fixed-basket log-price response to cross-market hospital mergers at specified horizons, retaining baseline patients and separating prices from quality and access. This attempt solves a simple bargaining model in which common insurer negotiations connect otherwise distinct patient markets. It also specifies what a distance-bin experiment identifies and why a universal geographical cutoff does not follow. No estimate for US mergers or current regulatory threshold is asserted.

Consider hospital systems $A$ and $B$ operating in separate patient markets with no direct patient substitution. A common insurer's gross value from a network is $V(S)$ for $S\subseteq\{A,B\}$, with $V(\varnothing)=0$, $V(A),V(B)>0$, and $V(AB)>\max\{V(A),V(B)\}$. These values summarize premiums and enrollment in a stipulated pricing environment. Hospital treatment resource costs are normalized to zero solely for the bargaining calculation. Transfer payments are $p_A,p_B$, and hospital bargaining weight is $\alpha\in(0,1)$. To preserve insurer participation relative to rejecting both agreements, also require $\alpha[2V(AB)-V(A)-V(B)]\le V(AB)$; the numerical examples below can both use $\alpha=1/2$.

Before merger, each bilateral negotiation takes the other agreement as fixed. The Nash product for agreement with $A$ is $p_A^\alpha[V(AB)-V(B)-p_A]^{1-\alpha}$, and analogously for $B$. After merger the joint system bargains an all-or-nothing contract with payment $p_{AB}$ and surplus $V(AB)$ relative to no agreement.

\begin{proposition}[An exact cross-market payment comparison]
The unique interior bargaining payments are
\[
 p_A=\alpha[V(AB)-V(B)],\quad
 p_B=\alpha[V(AB)-V(A)],\quad
 p_{AB}=\alpha V(AB).
\]
The merger changes total payment by
\[
 p_{AB}-(p_A+p_B)=\alpha[V(A)+V(B)-V(AB)]. \tag{1}
\]
\end{proposition}
\begin{proof}
For a positive bilateral incremental surplus $D$, maximize $\alpha\log p+(1-\alpha)\log(D-p)$ on $0<p<D$. The derivative $\alpha/p-(1-\alpha)/(D-p)$ vanishes exactly at $p=\alpha D$, and the objective is strictly concave. Apply this calculation to each stipulated disagreement point and subtract. All incremental surpluses are positive under the assumptions, so logarithms and interior solutions are defined.
\end{proof}

If $V(A)=V(B)=10,V(AB)=15$, total payment rises from $10\alpha$ to $15\alpha$ despite the absence of patient substitution. If instead $V(AB)=25$, payment falls from $30\alpha$ to $25\alpha$. The sign depends on the insurer's portfolio value and disagreement points. This is not a universal theory of cross-market merger effects: different bargaining protocols, outside networks or capacity constraints change those primitives.

A single negotiated basket can be defined with one unit from each hospital, so its premerger price is $p_A+p_B$ and postmerger price $p_{AB}$. The positive-price log effect in the first example is $\log(3/2)$. It is an accounting consequence of the chosen transfers, not a patient welfare calculation. There is no claim that commercial basket prices in actual negotiations are determined by this one-contract model.

\section{A concrete role for network overlap}
Parameterize $V(AB)=V(A)+V(B)-\omega K$, where $K\ge0$ measures a common-insurer or employer connection and $\omega\ge0$ is a bargaining-environment parameter. Maintain the positive incremental surplus restrictions. Formula (1) then gives a merger payment change $\alpha\omega K$. Distance is absent unless a further relation specifies how $K$ depends on distance.

Two mergers at the same distance can therefore have different effects through different $K$. Conversely, two very distant systems can have the same effect as nearby systems if the same insurer connection links them. A distance cutoff that presumes $K=0$ beyond a threshold is an additional economic restriction. Merely observing different hospital referral regions does not prove that insurer disagreement values separate.

The example also warns against mechanically controlling for postmerger network overlap. If merger itself changes network contracts, that overlap is a mediator; conditioning on it can remove part of the total merger effect. A useful heterogeneity variable is the preannouncement network relationship, measured before merger choices or anticipation affect it.

\section{An identifiable bin-average benchmark}
Let $B_d$ and $B_K$ be prespecified positive-probability bins of premerger distance and network overlap. Define a baseline merger-candidate population and a fixed cohort-weighted log basket price $Y_m(a)$ at the chosen horizon. Suppose completion $D_m$ is randomized within these bins, with probability $e_m\in(0,1)$ known conditional on baseline strata. Then
\[
 \tau(B_d,B_K)=\E[Y(1)-Y(0)\mid d\in B_d,K\in B_K]
\]
is identified by the corresponding conditional inverse-probability contrast
\[
 \E\left[\frac{DY}{e_m}-\frac{(1-D)Y}{1-e_m}
       \,\middle|\,d\in B_d,K\in B_K\right]. \tag{2}
\]
\begin{proof}
Condition on the baseline variables and the full potential-price pair. Independent assignment makes the expected weighted treated observation equal to $Y(1)$ and the expected weighted untreated observation equal to $Y(0)$. Subtract and integrate over the declared bin and fixed population weights.
\end{proof}

An observational difference-in-differences comparison instead needs its own conditional parallel-trend restriction. Such a restriction commonly identifies an effect for completed mergers within the compared strata, not automatically the effect for all candidate mergers in (2). Transport from the completed-merger population requires an additional response-selection restriction. Concurrent transactions must either be held fixed in the intervention or represented by a complete exposure vector; using affected controls violates the simple two-potential-price notation.

For continuous distance, a mean at an isolated point is only an almost-everywhere version without continuity. Suppose a known $L$-Lipschitz profile is imposed and a bin $B$ has width $h$ containing distance $d$. Then
\[
 |\tau(d)-\E[\tau(D)\mid D\in B]|\le Lh,
\]
by integrating $|\tau(d)-\tau(D)|\le L|d-D|\le Lh$. This deterministic interpolation error must accompany the sampling interval for the bin average. More observations reduce sampling error; they do not justify an unsupported $L$ or extrapolation to unobserved distances.

\section{Price increases and welfare remain separate}
With unchanged treatment resources and health, a higher negotiated payment transfers income from insured households or insurer owners to hospital owners. Equal dollar welfare weights make that pure transfer cancel. If patient health changes by $\Delta Q$, real resources by $\Delta C$, and different ownership weights apply, welfare requires those separate terms and the pass-through of payments into premiums. The sign of (1) alone is insufficient.

Similarly, apparent prices can change because of coding, service intensity, risk composition or basket substitution. Fixing the basket and patient mix before merger makes the estimand explicit but does not guarantee those data can be measured. The unresolved task is a credible merger counterfactual, validated premerger network exposure, informative distance support and separate quality/access outcomes. The bargaining theorem is a mechanism possibility and the bin formula an identification benchmark, not an empirical merger assessment.

\begin{thebibliography}{9}
\bibitem{rand} J. L. Liu, Z. M. Levinson, A. Zhou, X. Zhao, P. Nguyen and N. Qureshi (2022), \emph{Environmental Scan on Consolidation Trends and Impacts in Health Care Markets}, RAND RR-A1820-1. \url{https://www.rand.org/pubs/research_reports/RRA1820-1.html}. The checked primary report page discusses consolidation prices, quality and access. It is not cited as a proof of a geographical cutoff or of the simplified bargaining equations.
\end{thebibliography}

\end{document}
